\subsection{有限迹自同态}

我们回顾，E 的正自同态 \(u\) 是有限迹的，当且仅当存在 E 的标准正交基 \((e_i)_{i \in I}\) ，使得

\[
\sum_ {i \in \mathrm{I}} \langle e _ {i} | u (e _ {i}) \rangle <   + \infty
\]

（EVT, V, p. 48，引理 3 及 p. 49，定义 7）。若 K = C，E 的有限迹自同态所成的空间 \(\mathcal{L}_1(\mathrm{E})\) 是由正的有限迹自同态组成的集生成的向量空间（EVT, V, p. 50，定义 8）；若 K = R，空间 \(\mathcal{L}_1(\mathrm{E})\) 定义为交 \(\mathcal{L}(\mathrm{E}) \cap \mathcal{L}_1(\mathrm{E}_{(\mathbf{C})})\) （EVT, V, p. 50）。

若 \(u\in \mathcal{L}_1(\mathrm{E})\) ，则级数

\[
\sum_ {i \in \mathrm{I}} \langle e _ {i} | u (e _ {i}) \rangle
\]

对 E 的每个标准正交基 \((e_{i})_{i\in\mathrm{I}}\) 收敛，且其和与标准正交基无关；这个和称为 u 的迹 \(\operatorname{Tr}(u)\) （EVT, V, p. 50）。若 K = R，我们有 \(\operatorname{Tr}(u) = \operatorname{Tr}(u_{(\mathbf{C})})\) 。

设 \(u \in \mathcal{L}_1(\mathrm{E})\) 。我们有 \(u^* \in \mathcal{L}_1(\mathrm{E})\) ，且 \(\operatorname{Tr}(u^*) = \overline{\operatorname{Tr}(u)}\) （同前）。

命题 12. --- 设 \(B = (e_i)_{i \in I}\) 是 E 的标准正交基。设 \(u \in \mathcal{D}_B(E)\) ，并用 \(\lambda = (\lambda_i)_{i \in I}\) 表示它的特征值族。自同态 \(u\) 是有限迹的，当且仅当族 \(\lambda\) 在 K 中可和。此时我们有 \(\operatorname{Tr}(u) = \sum \lambda_i\) 。

如有必要，把 \(u\) 换成 \(u_{(\mathbf{C})}\) ，我们可设 \(K = \mathbf{C}\) 。

先设 \(u\) 是有限迹的。根据 EVT, V, p. 50 及 p. 49 的公式 (25)，族 \((\langle e_i | u(e_i) \rangle)_{i \in I} = (\lambda_i)_{i \in I}\) 是可和的。

反之，设族 \(\lambda\) 是可和的。

于是族 \((\mathcal{R}(\lambda_{i})^{+})\) 、\((\mathcal{R}(\lambda_{i})^{-})\) 、\((\mathcal{I}(\lambda_{i})^{+})\) 、\((\mathcal{I}(\lambda_{i})^{-})\) 都是可和的。自同态 u 是 \(\mathcal{D}_{\mathrm{B}}(\mathrm{E})\) 中以这些族为特征值的元素的线性组合。\(\mathcal{D}_{\mathrm{B}}(\mathrm{E})\) 的这些元素都是正的。由于按定义，有限迹自同态空间由正的有限迹自同态生成，因此可设对所有 i 有 \(\lambda_{i} \geqslant 0\) 。由于 \(\langle e_{i} | u(e_{i}) \rangle = \lambda_{i}\)，族 \((\langle e_{i} | u(e_{i}) \rangle)_{i \in \mathrm{I}}\) 是可和的，故 u 是有限迹的（EVT, V, p. 48，引理 3）。

最后，若 u 有有限迹，则由 EVT, V, p. 50，我们有

\[
\operatorname{Tr} (u) = \sum_ {i \in \mathrm{I}} \left\langle e _ {i} \mid u \left(e _ {i}\right) \right\rangle = \sum_ {i \in \mathrm{I}} \lambda_ {i}.
\]

推论 1. --- 设 u 是 E 的有限迹自同态。

\begin{enumerate}
\def\labelenumi{\alph{enumi})}
\item
  自同态 u 是紧的；
\item
  设 \(\mathrm{B} = (e_i)_{i \in \mathrm{I}}\) 是 u 的始空间 \(\operatorname{Ker}(u)^{\circ}\) 的标准正交基，\(\mathrm{C} = (f_i)_{i \in \mathrm{I}}\) 是 E 的标准正交族，且 \((\alpha_i)_{i \in \mathrm{I}}\) 是 \((\mathbf{R}_+^*)^{\mathrm{I}}\) 中的族，使得 \(u(e_i) = \alpha_i f_i\) 对所有 \(i \in I\) 成立（IV, p. 150 的推论 2）。我们有
\end{enumerate}

\[
\operatorname{Tr} (u) = \sum_ {i \in I} \alpha_ {i} \langle e _ {i} | f _ {i} \rangle .
\]

为证明 \(a\)），可设 \(\mathbf{K} = \mathbf{C}\) （EVT, V, p. 50 及 III, p. 2 的注 4），且 \(u\) 是正的（EVT, V, p. 50，定义 8）。

由 EVT, V, p. 56 的推论 1，存在 E 的标准正交基 \(B = (e_i)_{i \in I}\) ，使得 \(u\) 在基 B 中是对角的，且特征值族 \(\lambda\) （\(u\) 的特征值族）属于 \(\mathbf{R}_+^I\) 并可和。因此族 \(\lambda\) 属于 \(\mathcal{C}_0(\mathrm{I};\mathrm{K})\) （TG, III, p. 38，命题 1），从而自同态 \(u\) 是紧的（IV, p. 148 的命题 2）。

我们来证明 \(b\)）。设 \((e_i)_{i \in \mathbf{J}}\) 是 E 的标准正交基，它扩张族 B，使得 \(u(e_i) = 0\) （当 \(i \in \mathbf{J} - \mathbf{I}\) 时）。我们得到

\[
\operatorname{Tr} (u) = \sum_ {i \in \mathrm{J}} \langle e _ {i} \mid u (e _ {i}) \rangle = \sum_ {i \in \mathrm{I}} \langle e _ {i} \mid u (e _ {i}) \rangle = \sum_ {i \in \mathrm{I}} \alpha_ {i} \langle e _ {i} \mid f _ {i} \rangle .
\]

推论 2. --- 设 F 是希尔伯特空间。设 \(u \in \mathcal{L}(\mathrm{E};\mathrm{F})\) 是希尔伯特--施密特算子。则 u 是紧线性算子。

设 \((j,|u|)\) 是 u 的极分解（I, p. 140，定义 4）。按定义（EVT, V, p. 50，定义 9），自同态 \(u^{*}u\) 有有限迹，因而是紧的（推论 1, a)）；对 \(|u|=\sqrt{u^{*}u}\) （III, p. 91 的命题 6, b)）与 \(u=j|u|\) （III, p. 5 的命题 3）结论同样成立。

推论 3. --- 假设 K = C。设 u 是 E 的正紧自同态。自同态 u 是有限迹的，当且仅当它的特征值的递减族 \((\lambda_{n}(u))_{n\in\mathrm{I}_{\mathrm{E}}}\) 可和；此时 u 的迹就是这个族的和。

由 IV, p. 149 的定理 1，这可由前一命题以及序列 \((\lambda_n(u))_{n\in \mathrm{I_E}}\) 的定义（IV, p. 151 的 no. 3）得出，并用到 EVT, V, p. 49 的公式 (28)。

推论 4. --- 设 F 是希尔伯特空间。设 \(u \in \mathcal{L}(\mathrm{E};\mathrm{F})\) 是紧线性映射。设 \(\mathrm{B} = (e_i)_{i \in \mathrm{I}}\) 是始空间 \(\operatorname{Ker}(u)^{\circ}\) （\(u\) 的始空间）的标准正交基，\(\mathrm{C} = (f_i)_{i \in \mathrm{I}}\) 是 F 的标准正交族，且 \((\alpha_{i})_{i\in\mathrm{I}}\) 是 \((\mathbf{R}_{+}^{*})^{\mathrm{I}}\) 中的族，使得 \(u(e_{i})=\alpha_{i}f_{i}\) 对所有 \(i\in I\) 成立（IV, p. 150 的推论 2）。

E 的自同态 \(|u|\) 有有限迹，当且仅当族 \((\alpha_{i})_{i\in \mathrm{I}}\) 可和。此时我们有 \(u\in \mathcal{L}_2(\mathrm{E};\mathrm{F})\) ，并且

\[
\operatorname{Tr} (| u |) = \sum_ {i \in \mathrm{I}} \alpha_ {i}, \quad \| u \| _ {2} ^ {2} = \operatorname{Tr} (u ^ {*} u) = \sum_ {i \in \mathrm{I}} \alpha_ {i} ^ {2},\tag{8}
\]

特别地有 \(\|u\|_{2}\leqslant\mathrm{Tr}(|u|)\) 。

\(|u|\) 的非零特征值族是 \((\alpha_{i})_{i\in \mathrm{I}}\) ，故 \(|u|\) 有有限迹当且仅当族 \((\alpha_{i})_{i\in \mathrm{I}}\) 可和（命题 12）。若确实如此，非零特征值族 \((\alpha_{i}^{2})\) （\(u^{*}u = |u|^{2}\) 的非零特征值族）可和，且公式 (8) 由同前得出。

引理 7. --- 设 \(\lambda = (\lambda_i)_{i \in \mathrm{I}}\) 是复数族。对每个 \(t \in \mathbf{C}^*\) ，用 \(n_t\) 表示 \(i \in \mathrm{I}\) 中使 \(\lambda_i = t\) 者所成之集的基数。族 \((\lambda_i)_{i \in \mathrm{I}}\) 可和，当且仅当 \(n_t\) 对所有 \(t \in \mathbf{C}^*\) 有限，且族 \((n_t t)_{t \in \mathbf{C}^*}\) 可和。此时这两个族的和相等。

设族 \((\lambda_{i})_{i\in\mathrm{I}}\) 可和。对每个 \(t\in\mathbf{C}\) ，用 \(I_{t}\) 表示 \(i\in I\) 中使 \(\lambda_{i}=t\) 者所成之集。把 TG, III, p. 39 的定理 2 应用于 I 被诸集 \(I_{t}\) 所作的分划，可知集 \(I_{t}\) 对每个 \(t\in\mathbf{C}^{*}\) 有限，且族 \((n_{t}t)_{t\in\mathbf{C}^{*}}\) 可和，其和等于族 \((\lambda_{i})_{i\in\mathrm{I}}\) 的和。

反之，设 \(n_{t}\) 对所有 \(t \in \mathbf{C}^{*}\) 有限，且族 \((n_{t}t)_{t \in \mathbf{C}^{*}}\) 可和。设 J 是 I 的有限子集，\(\Lambda\) 是诸 \(\lambda_{i}\) （\(i \in J\)）所成之集。我们有

\[
\left| \sum_ {i \in J} \lambda_ {i} \right| \leqslant \sum_ {t \in \Lambda - \{0 \}} n _ {t} | t | \leqslant \sum_ {t \in \mathbf {C} ^ {*}} n _ {t} | t |,
\]

因此族 \((\lambda_{i})_{i\in\mathrm{I}}\) 可和（TG, VII, p. 17，推论）。

命题 13. --- 假设 K = C。设 u 是 E 的紧正规自同态。对 \(t \in \mathrm{Sp}_{\mathrm{s}}(u)\) ，用 \(n_{t} \geqslant 1\) 表示 u 的特征值 t 的谱重数。u 有有限迹的必要充分条件是族 \((n_{t}t)_{t \in \mathrm{Sp}_{\mathrm{s}}(u)}\) 可和。此时 u 的迹就是这个族的和。

设 \(\mathrm{B} = (e_i)_{i \in \mathrm{I}}\) 是 E 的标准正交基，使得 u 在基 B 中是对角的（IV, p. 149 的定理 1），并设 \(\lambda = (\lambda_i)_{i \in \mathrm{I}}\) 是它的特征值族。由于对 \(t \in \mathrm{Sp}_\mathrm{s}(u)\) ，重数 \(n_t\) 等于元素 \(i \in I\) 中使 \(\lambda_i = t\) 者的个数，且 \(\lambda\) 的非零元素属于 \(\mathrm{Sp}_{\mathrm{s}}(u)\) ，断言即由命题 12 与前一引理得出。

